\section{Introduction}
\label{sec:intro}

In 1961, the Atlas computer introduced demand paging: programs could
address more memory than physically existed, and the hardware would
transparently load pages from backing store as
needed~\citep{kilburn1962one}. Before Atlas, programmers manually
managed overlays---explicitly swapping code segments in and out of a
fixed address space. The transition from overlays to virtual memory
is one of the foundational advances in systems architecture.

Large language model context windows are in the overlay era. Every
agentic AI tool---Claude Code, Cursor, GitHub Copilot,
Windsurf---manually assembles context on every API call: tool
definitions, system prompts, and the complete message history. There
is no eviction policy, no demand loading, no working set estimation.
The default behavior is accumulation until the context limit forces
a crisis.

This paper makes two contributions. First, we provide the empirical
evidence that context management is a memory management problem. We
instrument 857~production sessions (54{,}170 API calls,
4.45~billion effective input tokens) and find that 21.8\% of tokens
are structural waste from three sources: unused tool schemas
(11.0\%), duplicated content (2.2\%), and stale tool results (8.7\%)
reprocessed at a median amplification of 84.4$\times$. These are
not implementation bugs. They are the inevitable consequence of
managing a finite resource (the context window) without the
abstractions that operating systems developed for the analogous
problem fifty years ago.

Second, we build the system. \textsc{Pichay} is a demand paging
system for LLM context windows, implemented as a transparent proxy
between client and inference API. It distinguishes garbage collection
(removing ephemeral tool output that cannot be re-requested) from
paging (evicting addressable content that can be faulted back in).
It detects page faults when the model re-requests evicted content.
It pins working-set pages identified by fault history. And the
eviction summaries it generates---designed as space-saving
markers---function as retrieval handles that models understand
without instruction: ``\texttt{[Paged out: Read file.py (8,192
bytes). Re-read if needed.]}''

In production deployment, the system extends a session from 7\%
free context to 43\% free---the difference between imminent context
death and substantial remaining capacity. Over a 681-turn session,
it sustains operation at 97\% eviction rate, though this extreme
case reveals the thrashing pathology well-known from OS
research~\citep{denning1968working}: when the working set exceeds
the resident set, the system spends more effort on faulting than on
useful work.

The deeper observation is that the context window is not memory. It
is L1~cache. And nobody is building the rest of the hierarchy.
The field's response to context limits is to make L1 bigger---1M
token windows, 10M token windows. This is the equivalent of building
machines with more physical RAM instead of inventing virtual memory.
It works, but it doesn't scale, and it misses the architectural
insight: a managed hierarchy of small-fast-expensive and
large-slow-cheap storage, connected by eviction policies and fault
mechanisms, outperforms any single level no matter how large.

We describe the full hierarchy: L1~(the generation window), L2~(the
working set, demand-paged and pinned), L3~(session history,
compressed with declared losses), and L4~(cross-session persistent
memory, indexed and retrievable). The first two levels are deployed
and evaluated. L3~is implemented---the model can compress conversation
turns into outcome summaries via cooperative cleanup tags---but not
yet evaluated at scale. L4~is designed and its mechanisms identified.

The contributions:
\begin{enumerate}
  \item A large-scale empirical characterization of context
    window utilization in a production agentic AI corpus, establishing that
    21.8\% of tokens are structural waste with a measured taxonomy.
  \item \textsc{Pichay}, a demand paging system for LLM context
    windows with empirically measured fault rate (0.0254\% offline,
    deployed in production).
  \item Fault-driven pinning: a page replacement policy that
    learns from its own mistakes---one fault pins a page for the
    session, reducing repeated faults in steady-state workloads.
  \item Cooperative memory management via two side channels
    (phantom tools and cleanup tags): mechanisms enabling the model
    to voluntarily release cold pages, request cached faults, and
    compress conversation history---a new point in the design space
    not available in hardware memory hierarchies, where applications
    are non-cooperative.
  \item The architectural observation that LLM context management
    maps structurally (not merely metaphorically) to virtual memory,
    and that the full OS memory hierarchy---from cache replacement
    policies to working set theory---applies directly.
  \item Design of a four-level memory hierarchy for LLM systems,
    with the first three levels deployed (L1~eviction, L2~pinning,
    L3~cooperative compaction) and the first two evaluated.
\end{enumerate}

% ====================================================================
